% !TeX root = ../main.tex
\section{Locality Sweep}\label{sec:locality}
\rubric{Locality and memory-bound analysis}{20}{Working-set transitions identified, and the compression
of SIMD gains in the memory-bound regime explained.}

\begin{hypothesis}
  \todo{At which $N$ each cache transition should appear for this kernel, and how scalar and SIMD
  builds should diverge or converge beyond LLC.}
\end{hypothesis}

\begin{knobs}
  Fixed       & kernel \tbd; float32; unit stride; aligned; pinned to CPU \tbd \\
  Varied      & $N$ (dense grid across L1/L2/LLC/DRAM); variant \\
  Counters    & \todo{perf events recorded (\code{PERF=1}), e.g.\ LLC-load-misses} \\
  Script      & \code{SIMD_Advantage_Profiling/scripts/run.sh locality} \\
\end{knobs}

\begin{figure}[tbp]
  \centering
  \IfDataExists{locality}{%
    \begin{tikzpicture}
      \begin{axis}[acs, xmode=log, log basis x=2, ymode=log, xlabel={$N$ (elements)}, ylabel={Time per element (ns)},
                   legend pos=north west]
        \addplot+[iqr bars] table[x=n, y=ns_per_elem_median, iqr of=ns_per_elem, only if={variant}{scalar}] {data/locality.csv};
        \addlegendentry{scalar}
        \addplot+[iqr bars] table[x=n, y=ns_per_elem_median, iqr of=ns_per_elem, only if={variant}{autovec}] {data/locality.csv};
        \addlegendentry{autovec}
        \addplot+[iqr bars] table[x=n, y=ns_per_elem_median, iqr of=ns_per_elem, only if={variant}{native}] {data/locality.csv};
        \addlegendentry{native}
        % TODO: mark cache boundaries in elements: \acsvline{<N>}{L1d} \acsvline{<N>}{L2} \acsvline{<N>}{LLC}
      \end{axis}
    \end{tikzpicture}%
  }
  \caption{Time per element across cache levels. \todo{takeaway}}\label{fig:locality}
\end{figure}

\subsection{Analysis}
\todo{Where each transition occurs (bytes = arrays $\times$ element size $\times N$) and what the
counters show there; why the SIMD speedup compresses once memory traffic dominates.}
